\section{通用通行账户与通行等价日记}\label{sec:universal}

第一项任务是把通行的共同描述与个体赋值分离开来。即使已经指定运动发生在哪里、位于目标哪一侧、方向如何以及重复多少次，路径的效用仍未确定。通用定理识别出所有允许的赋值都必须经过的可加信息空间；纤维定理随后精确刻画：何时两份有序日记会向这一赋值类中的每一个赋值提供完全相同的输入。

\subsection{通用可加载体}

对每个 $m\in\M$ 以及 $0\le a<b\le R_m$，引入一个形式生成元 $e_m(a,b)$。令 $\mathbb Z^{(\mathcal E)}$ 为这些生成元上的自由阿贝尔群，并对下列细分关系取商：
\[
 e_m(a,c)=e_m(a,b)+e_m(b,c),
 \qquad 0\le a<b<c\le R_m.
\]
所得群记为 $\mathfrak L$。在路径的转折点和目标跨越点处切分，可定义
\[
 q:\Pth\to\mathfrak L
\]
即把相应生成元的等价类相加。细分关系保证 $q$ 与所选切分方式无关。

弱时间变换是与一个从 $[0,1]$ 到自身的连续、非减、满射映射作复合。它可以改变速度并插入停留，但不能颠倒访问状态的先后顺序。

对于阿贝尔群 $G$，令 $\operatorname{Val}_G(\Pth)$ 表示所有映射 $\Phi:\Pth\to G$ 的集合，这些映射对弱时间变换不变，并且对每一次端点相容的拼接可加。

\begin{theorem}[可加赋值的通用通行账户]\label{thm:universal}
令 $G$ 为任意阿贝尔群。映射 $\Phi:\Pth\to G$ 对弱时间变换不变且对端点相容拼接可加，当且仅当存在唯一的同态 $\bar\Phi:\mathfrak L\to G$ 使得
\[
 \Phi=\bar\Phi\circ q.
\]
等价地，
\[
 \operatorname{Hom}_{\mathbf{Ab}}(\mathfrak L,G)
 \cong
 \operatorname{Val}_G(\Pth).
\]
此外，$\mathfrak L$ 与四模块整数阶梯账本群典范同构，并且在该同构下 $q$ 恰好变为 $\ell$。
\end{theorem}

由于该定理本身不含具体赋值，它同时允许个体特定和模块特定的权重。对于个体 $i$，任何允许的通行赋值都具有形式
\[
 \Phi_i(\gamma)=\bar\Phi_i\bigl(q(\gamma)\bigr),
\]
其中 $\bar\Phi_i$ 是个体特定同态。共同的是账户 $q(\gamma)$；不同的是对该账户的赋值。

常值路径的零通行价值由 $c_x*c_x=c_x$ 与可加性推出。状态水平、持续时间、健康、消费等非通行后果在条件通行排序中保持固定。

定理~\ref{thm:universal} 使用完整定义域 $\Pth$，其中包含每个典范生成元以及细分所需的所有拼接。受限菜单则需要各自的丰富性条件。

\subsection{通行账户的精确纤维}

令 $\approx_p$ 为下列操作生成的等价关系：
\begin{enumerate}
 \item 弱时间变换，包括插入或删除停留；
 \item 把所有常值路径认作零通行类；
 \item 交换以同一状态为基点的两个闭合游程；
 \item 对闭合路径作循环旋转，即对同一闭合日记选择不同的记录起点。
\end{enumerate}
每一种操作都保持有向账户不变。逆命题更微妙：需要证明账本系数相同的两条路径可以通过一系列操作连接，并且每个中间对象仍是可行的单一路径，同时目标达成标签始终保持不变。

\begin{theorem}[通行等价的恢复日记]\label{thm:fiber}
对于固定目标区间上的有限步路径，
\[
 q(\gamma)=q(\eta)
 \quad\Longleftrightarrow\quad
 \gamma\approx_p\eta.
\]
\end{theorem}

证明分三步使用几何结构。首先，把有限多条连续路径约化为同一个带标签路径图上的顶点词；其次，利用流量不平衡恢复每个开放实现的有序端点；最后，通过带根的叶节点删除，把每个实现写成一条约化核心路径加上一组闭合叶游程。闭合游程交换和插入点不变性使这些多重集能够对齐。只有对非零闭合账本选择共同基点时才需要循环旋转。若删除的叶节点就是目标，则恢复的游程恰好包含一个以目标为终点的趋向目标 单元和一个从目标出发的远离目标 单元，因此目标达成约定得到保留。完整论证见附录~\ref{app:fiber}。图~\ref{fig:fiber} 展示同基点闭合游程交换这一生成操作。

\begin{figure}[pos=htbp,align=\centering]
 \centering
 \includegraphics[width=0.98\textwidth]{Figure_2_path_fiber.pdf}
 \caption{\textbf{路径纤维内部的闭合游程交换。} 两份有序日记共享同一开放骨架，并包含以目标为共同基点的两个相同闭合游程，但先后顺序相反。交换这些同基点游程会保持每一个穿孔有向剖面以及各侧的目标达成次数。因此 $q(\gamma)=q(\eta)$ 且 $\gamma\approx_p\eta$。该图展示了路径纤维等价关系的一个生成操作。}
 \label{fig:fiber}
\end{figure}

每一个可加通行赋值都会给账本等价的历史赋予相同通行价值。上述操作穷尽了有向通行账户所丢弃的信息；在更大的应用模型中，非通行后果仍可能区分这些历史。

该结果利用有限的一维路径几何。经典的片段反转与分离--吸收结果允许更广泛的变换 \citep{Abrham1980,FleischnerEtAl1992}；本文定理则只使用始终停留在实际路径域中的账户保持操作，给出精确的纤维描述。

\subsection{由总通行与端点得到的一维压缩}\label{sec:compression}

对于侧 $s\in\{L,R\}$ 和穿孔距离 $r>0$，令 $N_s^T(r;\gamma)$ 与 $N_s^A(r;\gamma)$ 分别表示 趋向目标 与 远离目标 的跨越次数，并定义
\[
 D_s^\gamma=N_s^T-N_s^A,
 \qquad
 V_s^\gamma=N_s^T+N_s^A.
\]
当 $x$ 位于侧 $s$ 且距目标至少为 $r$ 时，令 $\chi_s(x,r)=1$；否则令其为零。

\begin{lemma}[净通行守恒]\label{lem:net}
对任意有限步路径和任意穿孔距离水平，
\[
 D_s^\gamma(r)
 =\chi_s(\gamma(0),r)-\chi_s(\gamma(1),r).
\]
因此，净有向通行完全由有序端点决定。
\end{lemma}

\begin{proposition}[由总通行与有序端点无损重构]\label{prop:gross-endpoints}
在有限步一维路径域上，记录
\[
 \bigl(V_L^\gamma,V_R^\gamma,\gamma(0),\gamma(1)\bigr)
\]
即可重构完整的有向通行账户 $q(\gamma)$。因此，在 BPCC 下，该记录相同意味着原始偏好上的无差异。
\end{proposition}

\begin{proof}
引理~\ref{lem:net} 由有序端点恢复 $D_s$，因此
\[
 N_s^T=(V_s+D_s)/2,
 \qquad
 N_s^A=(V_s-D_s)/2.
\]
对于有限步路径，$N_s^T$ 在零点附近足够小的穿孔邻域上为常数。其稳定的近目标值恰好等于实际从侧 $s$ 到达目标的趋向目标 通行次数。因此，穿孔有向剖面与目标轨迹都得到恢复，这正是完整账本。
\end{proof}

在严格方向单调性下，仅有总通行并不充分：同一穿孔区间上的典范 趋向目标 通行与其反向 远离目标 通行具有相同的总剖面，却分别位于停留的不同偏好方向。有序端点提供了缺失的方向不平衡。这一重构结果依赖于一维有限步几何。
